\section*{AI工具使用声明}
本参赛队在竞赛过程中使用了AI工具，主要用于题意分析、文献线索整理、程序设计与核验、图表制作和文稿编辑，详细使用情况见支撑材料《AI工具使用详情》。

所用工具为OpenAI Codex。本轮由AI将公共模型与各问求解分章整理，编辑摘要和论证文字，核对既有结果并编译检查；未重新计算模型，核心方程和数值结论保留。参赛队最终人工审阅尚待完成，使用者须依据真实记录确认披露，独立理解和核对方程、程序、结果与引用，并对提交内容负责。

\begin{thebibliography}{9}
\small\raggedright
\bibitem{silva2014turmeric}
da Silva W. P., e Silva C. M. D. P. S., Gomes J. P.
Comparison of boundary conditions to describe drying of turmeric
(Curcuma longa) rhizomes using diffusion models[J].
Journal of Food Science and Technology, 2014, 51(11): 3181--3189.
DOI: \href{https://doi.org/10.1007/s13197-012-0813-x}{10.1007/s13197-012-0813-x}.

\bibitem{silva2014coupling}
da Silva W. P., e Silva C. M. D. P. S., Gama F. J. A.
Estimation of thermo-physical properties of products with cylindrical
shape during drying: The coupling between mass and heat[J].
Journal of Food Engineering, 2014, 141: 65--73.
DOI: \href{https://doi.org/10.1016/j.jfoodeng.2014.05.010}{10.1016/j.jfoodeng.2014.05.010}.

\bibitem{tzempelikos2015quince}
Tzempelikos D. A., Mitrakos D., Vouros A. P., Bardakas A. V.,
Filios A. E., Margaris D. P.
Numerical modeling of heat and mass transfer during convective
drying of cylindrical quince slices[J].
Journal of Food Engineering, 2015, 156: 10--21.
DOI: \href{https://doi.org/10.1016/j.jfoodeng.2015.01.017}{10.1016/j.jfoodeng.2015.01.017}.

\bibitem{seyedabadi2017ale}
Seyedabadi E., Khojastehpour M., Abbaspour-Fard M. H.
Convective drying simulation of banana slabs considering non-isotropic
shrinkage using FEM with the Arbitrary Lagrangian--Eulerian method[J].
International Journal of Food Properties, 2017, 20(sup1): S36--S49.
DOI: \href{https://doi.org/10.1080/10942912.2017.1288134}{10.1080/10942912.2017.1288134}.

\bibitem{adrover2020pears}
Adrover A., Venditti C., Brasiello A.
A Non-Isothermal Moving-Boundary Model for Continuous and
Intermittent Drying of Pears[J].
Foods, 2020, 9(11): 1577.
DOI: \href{https://doi.org/10.3390/foods9111577}{10.3390/foods9111577}.
\end{thebibliography}
